% Faithful translation baseline; Chinese source PDF pp. 42--47.
\section{Schr\"odinger Equation in the Coordinate Representation and the WKB Approximation}
Consider the abstract state $\ket{\psi,t}$, whose Schr\"odinger equation is
\[\ii\hbar\frac{\partial}{\partial t}\ket{\psi,t}=\op H\ket{\psi,t}.\]
Multiply from the left by the eigenbra $\bra{\vect x}$ of the position operator $\vect{\op x}$. Since these eigenstates are time independent in the Schr\"odinger picture,
\[\ii\hbar\frac{\partial}{\partial t}\braket{\vect x}{\psi,t}
=\mel{\vect x}{\op H}{\psi,t}
=\int\dd^3x'\,\mel{\vect x}{\op H}{\vect x'}\braket{\vect x'}{\psi,t}.\]
Let
\[\op H=\frac{\vect{\op p}^{\,2}}{2m}+V(\vect{\op x}).\]
Since
\[\mel{\vect x}{\op H}{\vect x'}
=\delta^{(3)}(\vect x-\vect x')
\left[-\frac{\hbar^2}{2m}\nabla^2+V(\vect x)\right],\]
we obtain
\[\ii\hbar\frac{\partial}{\partial t}\braket{\vect x}{\psi,t}
=\left[-\frac{\hbar^2}{2m}\nabla^2+V(\vect x)\right]\braket{\vect x}{\psi,t}.\]
Define the wavefunction $\psi(\vect x,t)=\braket{\vect x}{\psi,t}$. This is the Schr\"odinger equation in the coordinate representation.

\subsection*{1) Time-independent Schr\"odinger equation}
Put
\[\psi(\vect x,t)=\varphi(\vect x)T(t).\]
If $\op H$ has no time dependence, inserting this into the Schr\"odinger equation separates it into two equations. The time-dependent factor is
\[T(t)=e^{-\ii Et/\hbar}.\]
The energy $E$ and $\varphi(\vect x)$ are determined by the other equation, the energy eigenvalue equation:
\[\boxed{\left[-\frac{\hbar^2}{2m}\nabla^2+V(\vect x)\right]
\varphi_E(\vect x)=E\varphi_E(\vect x).}\]
A stationary state is
\[\psi_E(\vect x,t)=\varphi_E(\vect x)e^{-\ii Et/\hbar}.\]
For bound states the bound-state condition
\[\lim_{|\vect x|\to\infty}\varphi_E(\vect x)=0\]
leads to discrete values of $E$.

Stationary states are particular solutions when $\op H$ is time independent. The general solution is a linear superposition of stationary states:
\[\psi(\vect x,t)=\sum_Ec_E\psi_E(\vect x,t).\]

\subsection*{2) Probability interpretation of the wavefunction}
The probability density of finding the particle at a given point is
\[\rho(\vect x,t)=|\braket{\vect x}{\psi(t)}|^2
=|\psi(\vect x,t)|^2=\psi^*(\vect x,t)\psi(\vect x,t).\]
When $V(\vect x)$ is real, the Schr\"odinger equation can be written as the continuity equation,
\[\boxed{\frac{\partial\rho(\vect x,t)}{\partial t}
+\vect\nabla\cdot\vect j(\vect x,t)=0,}\]
\[\vect j(\vect x,t)=-\frac{\ii\hbar}{2m}
\left[\psi^*(\vect x,t)\vect\nabla\psi(\vect x,t)
-\psi(\vect x,t)\vect\nabla\psi^*(\vect x,t)\right].\]
What is the physical meaning of $\vect j$? Integrate the continuity equation over a finite volume $V$:
\[\int_V\dd^3r\,\frac{\partial\rho(\vect r,t)}{\partial t}
+\int_V\dd^3r\,\vect\nabla\cdot\vect j(\vect r,t)=0,\]
\[\frac{\dd}{\dd t}\int_V\dd^3r\,\rho(\vect r,t)
=-\oint_{\partial V}\dd\vect S\cdot\vect j(\vect r,t).\]
Local probability conservation says that the change in the probability inside $V$ equals the negative of the outward probability flux across its boundary $S$. Thus $\vect j$ is called the probability-current density.

If we integrate over all space,
\[\frac{\dd}{\dd t}\int_{\mathbb R^3}\dd^3r\,\rho(\vect r,t)
=-\oint_{\infty}\dd\vect S\cdot\vect j(\vect r,t)=0.\]
Therefore $\int\dd^3r\,\rho(\vect r,t)$ is independent of time and is a constant. This implies:
\begin{enumerate}[label=\alph*)]
\item Total probability is conserved; particles are neither created nor annihilated. This describes nonrelativistic quantum mechanics.
\item The normalization of a wavefunction is independent of time. Schr\"odinger evolution preserves normalization.
\item If $V$ is complex, probability is not conserved, and this can effectively describe creation and annihilation of particles.
\end{enumerate}

\subsection*{3) How does a quantum system approach a classical system?}
An operator in a general state has no definite value, whereas its expectation value does. Thus expectation values of quantum observables correspond to physical quantities in classical mechanics.

The classical equation of motion is
\[\frac{\dd O}{\dd t}=\{O,H\}_{\text{Poisson bracket}},\]
while the quantum equation of motion for the expectation value is
\[\frac{\dd\langle O\rangle}{\dd t}
=\frac{\langle[\op O,\op H]\rangle}{\ii\hbar}.\]
In the classical limit, observables commute, and $[\op O,\op H]\to0$. To recover the classical equation of motion, one requires
\[\hbar\longrightarrow0,\qquad
\frac{\langle[\op O,\op H]\rangle}{\ii\hbar}
\longrightarrow\{O,H\}_{\text{Poisson bracket}}.\]

\subsection*{4) WKB approximation to the time-independent equation}
Only a few Schr\"odinger equations can be solved exactly. Expanding the Schr\"odinger equation in powers of $\hbar$ and solving it by successive approximations gives the WKB approximation (Wentzel, Kramers and Brillouin).

Write a stationary-state wavefunction as
\[\Psi(\vect r,t)=\psi(\vect r)e^{-\ii Et/\hbar}
=e^{\ii s(\vect r)/\hbar}e^{-\ii Et/\hbar}
=e^{\ii[s(\vect r)-Et]/\hbar}.\]
Insert this into the stationary Schr\"odinger equation,
\[\left[-\frac{\hbar^2}{2m}\nabla^2+V(\vect r)\right]\psi(\vect r)
=E\psi(\vect r).\]
We find
\[\frac{(\vect\nabla s)^2}{2m}+V(\vect r)-E
-\frac{\ii\hbar}{2m}\nabla^2s=0.\]
Expand $s(\vect r)$ in $\hbar/\ii$:
\[s(\vect r)=s_0(\vect r)+\frac\hbar\ii s_1(\vect r)
+\left(\frac\hbar\ii\right)^2s_2(\vect r)+\cdots.\]
Equate like powers of $\hbar$:
\[\begin{array}{ll}
\hbar^0:&(\vect\nabla s_0)^2+2m[V(\vect r)-E]=0,\\[3pt]
\hbar^1:&\vect\nabla s_1\cdot\vect\nabla s_0+\frac12\nabla^2s_0=0,\\[3pt]
&\vdots
\end{array}\]
For one dimension,
\[s_0(x)=\pm\int^x p(x')\dd x'+C,\qquad
s_1(x)=-\ln\sqrt{p(x)}+\widetilde C,\]
\[\boxed{p(x)=\sqrt{2m[E-V(x)]}.}\]
Keeping terms through first order in $\hbar$,
\[\begin{split}
\psi_{\mathrm{II}}(x)
&=\exp\left[\frac\ii\hbar\left(s_0(x)+\frac\hbar\ii s_1(x)\right)\right]\\
&=\frac{C_1}{\sqrt{p(x)}}\exp\left(\frac\ii\hbar\int^xp(x')\dd x'\right)
+\frac{C_2}{\sqrt{p(x)}}\exp\left(-\frac\ii\hbar\int^xp(x')\dd x'\right).
\end{split}\]
This solution is valid where $p(x)=\sqrt{2m[E-V(x)]}>0$, that is, $E>V$. In classical mechanics, $E=T+V$; the condition $E>V$ therefore identifies a classically allowed region.

In the classically forbidden region, $E<V$, put
\[p(x)=\ii\widetilde p(x)=\ii\sqrt{2m[V(x)-E]}.\]
Then
\[\psi_{\mathrm{I,III}}(x)
=\frac{C_3}{\sqrt{|p(x)|}}\exp\left(\frac1\hbar\int^x\widetilde p(x')\dd x'\right)
+\frac{C_4}{\sqrt{|p(x)|}}\exp\left(-\frac1\hbar\int^x\widetilde p(x')\dd x'\right).\]
At $E=V$, the momentum $p=0$ and the solution diverges: WKB is no longer valid. A point $x_0$ satisfying $E-V(x_0)=0$ is a classical turning point. How can one solve the equation near $x_0$?

\begin{center}
\begin{tikzpicture}[>=Latex,scale=.9]
\draw[->](0,0)--(6.5,0)node[right]{$x$};\draw[->](0,0)--(0,3.0)node[above]{$V$};
\draw[thick,smooth] plot coordinates {(.6,2.5)(1.2,1.65)(2,1.25)(2.8,.9)(3.4,.8)(4.0,.95)(4.8,1.6)(5.6,2.7)};
\draw[dashed](.2,1.7)--(6.1,1.7)node[right]{$E$};
\fill (1.16,1.7) circle(2pt);\fill (4.89,1.7) circle(2pt);
\draw[dashed](1.16,0)--(1.16,1.7);\draw[dashed](4.89,0)--(4.89,1.7);
\node[below]at(1.16,0){$x_1$};\node[below]at(4.89,0){$x_2$};
\node at(.7,2.0){\scriptsize I};\node at(3,2.0){\scriptsize II};\node at(5.7,2.0){\scriptsize III};
\end{tikzpicture}
\end{center}
Expand the potential to linear order near $x_0$:
\[V(x)\simeq V(x_0)+\left.\frac{\dd V}{\dd x}\right|_{x_0}(x-x_0).\]
Substitution into the stationary Schr\"odinger equation gives
\[\frac{\dd^2\varphi}{\dd x^2}
-\frac{2m}{\hbar^2}\left.\frac{\dd V}{\dd x}\right|_{x_0}
(x-x_0)\varphi=0.\]
The exact solutions are Bessel functions of order $\pm\frac13$. Consider their asymptotic solutions when $|x-x_0|$ is large and require continuity with the classically allowed and forbidden solutions on either side of $x_0$. In this way the constants $C_1,C_2,C_3,C_4$ can be fixed.